% Generated draft: requires mathematical and source review.
\documentclass[11pt]{article}
\usepackage[margin=25mm]{geometry}
\usepackage{fontspec}
\setmainfont{Latin Modern Roman}
\usepackage{amsmath,amssymb,mathtools,unicode-math}
\setmathfont{Latin Modern Math}
\usepackage{enumitem,needspace,xurl,hyperref}
\setlength{\parindent}{0pt}
\setlength{\parskip}{5pt}
\setlength{\emergencystretch}{3em}
\newcommand{\N}{\mathbb{N}}
\newcommand{\Z}{\mathbb{Z}}
\newcommand{\Q}{\mathbb{Q}}
\newcommand{\R}{\mathbb{R}}
\newcommand{\C}{\mathbb{C}}
\newcommand{\F}{\mathbb{F}}
\newcommand{\A}{\mathbb{A}}
\newcommand{\PP}{\mathbb{P}}
\newcommand{\E}{\mathbb{E}}
\newcommand{\eps}{\varepsilon}
\newcommand{\dd}{\,\mathrm{d}}
\begin{document}
\section*{UnsolvedMath ID: 800875}
\textbf{Source problem code:} catalogue entry
\subsection{Definitions and mathematical statement}
Let \(n\geq 1\) be an integer and let \(\mathbb B=\{z\in\C^n:\|z\|<1\}\), where \(\|\cdot\|\) is the Euclidean norm. Let \(\varphi\) be plurisubharmonic on \(\mathbb B\), and assume that its complex Monge--Amp\`ere measure
\[
\mu_\varphi=(dd^c\varphi)^n
\]
is well defined. The source suggests local boundedness of \(\varphi\) on \(\mathbb B\setminus\{0\}\) as a possible hypothesis under which to consider this measure. It does not specify the construction or precise domain of definition of the Monge--Amp\`ere operator; that qualification remains unresolved here. No boundary condition on \(\partial\mathbb B\) is supplied.

The Lelong number of \(\varphi\) at the origin is
\[
\nu(\varphi,0)=\liminf_{z\to0,\ z\ne0}\frac{\varphi(z)}{\log\|z\|}.
\]
Saying that \(\mu_\varphi\) has a Dirac mass at the origin means that \(\mu_\varphi(\{0\})>0\). The source does not fix the normalization of \(dd^c\); multiplying it by a positive constant does not affect whether this mass is positive.

\textbf{Question (Guedj--Rashkovskii).} Under the preceding well-definedness assumption, is the implication
\[
(dd^c\varphi)^n(\{0\})>0
\quad\Longrightarrow\quad
\nu(\varphi,0)>0
\]
valid for every such \(\varphi\)? The suggested hypothesis of local boundedness away from the origin is retained as a source-specified setting, without identifying it with a fully specified operator domain.
\subsection{Short English statement}
Does a positive atom at the origin in the well-defined complex Monge--Amp\`ere measure of a plurisubharmonic function force a positive Lelong number there? The source suggests considering functions locally bounded away from the origin, but leaves the precise definition of the measure unspecified.
\subsection{Sources}
\begin{itemize}
\item Local UnsolvedMath registry: permanent ID \(800875\), exact source problem code \texttt{catalogue entry}; supplied complete source record.
\item Dinew, Guedj, Zeriahi, \emph{Open Problems in Pluripotential Theory} (2015), Question 6. Supplied source URL: \url{https://arxiv.org/abs/1511.00705}. The question itself is attributed in the statement to Guedj and Rashkovskii.
\item The supplied provenance records access on 2026-10-16 and extraction by \texttt{arxiv\_latex}. These are supplied provenance claims, not independent source verification. The registry labels the problem \texttt{open}; this is a supplied status label, not an independently verified literature assessment.
\end{itemize}
\end{document}
